\cotefile{137}
\scanpage{17}

Ces conditions, les entiers $\mu_i(i\geq1)$ $\nu_j$ $(j\geq0)$, $f_{i\ell}$ $(i\geq1,\ \ell\in\unsure{J|I|\mathfrak{B}})$ sont uniquement déterminés par les conditions suivantes

a) Pour \inserted{tout} $i\geq1$, les $f_{i\ell}\inserted{\ne0}$ sont les \inserted{multiplicités des} diviseurs élémentaires de $\deleted{\unsure{H_{i-1}|\mathrm{Tors}}}\pi_{i-1}(A_{\bullet})$ $(=H_{i-1}(L_{\bullet}^{!}))$

b) $\deleted{\unsure{k}}\unsure{k_i|\nu_i}=\operatorname{rg}\pi_i(A_{\bullet})$ $(=\operatorname{rg}H_i(L_{\bullet}^{!}))$

c) On a \deleted{rang} $\mu_i=\operatorname{rg}B_{i-1}(L_{\bullet}^{!})\deleted{\unsure{+f_i}}-f_i$ $(\text{ou } \unsure{\beta_i|f_i|\rho_i}=\sum_{\ell} f_{i,\ell})$

d) On a de plus

\[
\left\{\begin{aligned}
\operatorname{rg}Z_i\deleted{\unsure{(}}(L_{\bullet}^{!})&=\mu_{i+1}\deleted{\unsure{+k_i}}+f_{i+1}+\nu_i\\
\operatorname{rg}L_i^{!}&=\mu_{\unsure{i\ldotp|i+1}}+\deleted{\unsure{f_i+}}\unsure{f_{i+1}|\mu_{i+1}|\rho_{i+1}}+f_i+\nu_i .
\end{aligned}\right.
\]

Lorsque $L_{\bullet}$ est (str.) $\partial$ parfait \unsure{stt|sHt|ssi} le nb de facteurs non nuls qui interviennent est fini.

\uline{Dém} \deleted{\unsure{Le fonct.}} \inserted{On} la \unsure{théorie|topo|vertu} de Dold-Puppe, on réduit tout à une question de complexes de chaînes. Comme tt sous module d'un k-module projectif de t. f. de $k$ est itou, on conclut que \deleted{\unsure{tout}} tout complexe, s'il est \deleted{\unsure{de modules|pseudo}} \inserted{str.} à comp. proj. de t.f., s'écrit comme prod\textsuperscript{t} de complexes (définis à isom \deleted{\unsure{près}} \inserted{uniques} \deleted{\unsure{uniques}} près, \deleted{par le choix de} \inserted{-- mais l'isom} $\prod\simeq L_{\bullet}^{!}$ n'est pas unique) de la forme

\diagram{\begin{tikzcd}[ampersand replacement=\&,column sep=small,row sep=small]
K_i\mathpunct{:} \& {}-0 \arrow[r] \& \overset{\text{degré }i}{P_i} \arrow[r,hook] \& \overset{\text{degré }i-1}{Q_i} \arrow[r,dash] \& 0- \\
\& \& B_{i-1} \arrow[u,equal] \arrow[r,dash,hook] \& Z_{i-1} \arrow[u,equal] \&
\end{tikzcd}}

% archive sheet number: 16
\scanpage{18}

On désigne par $\overline P_i$\inserted{$\subset Q_i$} l'image inverse \inserted{dans $Q_i$} du sous-grpe de torsion de $Q_i/B_i\simeq H_{i-1}$, $N_{i-1}=Q_i/\overline P_i$ est de t.f. sans torsion, donc \margin{$H_{i-1}/\mathrm{torsion}$ $\nearrow$} projectif de t.f., donc $Q_i\simeq P_i\oplus N_{i-1}$ et le complexe $K_i$ est isom à la somme directe des complexes

\diagram{\begin{tikzcd}[ampersand replacement=\&,column sep=small,row sep=small]
\text{deg} \& \& \& \text{deg }i \& i-1 \& \& \\
\overline K_i \& \cdots\arrow[r] \& 0\arrow[r] \& P_i\arrow[r,hook] \& \overline P_i\arrow[r] \& 0\arrow[r] \& 0\\
N_i[-i] \& \& 0\arrow[r] \& 0\arrow[r,dash] \& N_{i-1}\arrow[r] \& 0 \&
\end{tikzcd}}

Enfin, $P_i\subset\overline P_i$ est une inclusion de modules projectifs de t.f. tels que $\overline P_i/P_i$ $(\simeq\mathrm{tors}\ H_{i-1})$ soit de torsion, et est justiciable de la théorie des diviseurs élémentaires : \deleted{On aura} \inserted{(si $k$ \uline{principal})} ce \deleted{$\unsure{\overline K_i\simeq|K_i\simeq}$} complexe est une somme directe de complexes de la forme

\diagram{\begin{tikzcd}[ampersand replacement=\&,column sep=small,row sep=0pt]
0\arrow[r,dash] \& \overset{i}{k}\arrow[r,"n_{i,\alpha}"] \& \overset{i-1}{k}\arrow[r] \& 0 \& 1\leq\alpha\leq N_i \\
\& \& \& \& \operatorname{rg}P_i=\operatorname{rg}\overline P_i \arrow[u,equal]
\end{tikzcd}}

avec les $n_{i\alpha}\in k-(0)$ ($n_{i\alpha}$ pouvant être une unité -- donc $1$ alors) ; le nombre des $n_{i\alpha}\in k^{*}$ qui intervient est noté $\mu_i$, (il lui correspond le facteur $(\Lambda^{i}\Phi)^{\otimes\mu_i}$) ; et si on prend les \inserted{\unsure{not.|not}} $n_{\unsure{j,a|i,\alpha}}$ tels qu'ils sont croissants pour la divisibilité, on trouve une suite unique ; les multiplicités \inserted{d'occurrence} des diviseurs \unsure{él.|é}

% archive sheet number: 17
\scanpage{19}

dans la suite des $\operatorname{div}(n_{i\alpha})$ est $f_{i\ell}$, et $\Phi(i,\ell)_k$ est le \deleted{complexe} \inserted{module $\frac12$ \unsure{simpl.|simplicial}} associé à

\diagram{\begin{tikzcd}[ampersand replacement=\&,column sep=small]
{}-{} \& 0\arrow[r] \& \overset{\mathrm{deg}\,i}{k}\arrow[r,"n"] \& \overset{\mathrm{deg}\,i-1}{k}\arrow[r] \& 0-
\end{tikzcd}}

où $\operatorname{div}(n)=\ell$. On a donc une extension

\diagram{\begin{tikzcd}[ampersand replacement=\&,column sep=small]
0\arrow[r] \& \Lambda^{i-1}\Psi\arrow[r] \& \Phi(i,n)\arrow[r,dash] \& \Lambda^i\Psi\arrow[r] \& 0
\end{tikzcd}}

où l'hom associé à cette extension

\[ k\longrightarrow k \]

(cf ---) est $n\,\mathrm{id}$, où $\operatorname{div}(n)=\ell$, donc

\[
\deleted{\unsure{H_j}}\pi_j(\Phi(i,n))=
\begin{cases}0&\text{si }j\ne i-1\\
k/\operatorname{div}n&\text{si }j=i-1.
\end{cases}
\]

\uline{Remarques} 1) La fonction de Hilbert de $\Phi(i,n)$ est donc la $=$ que celle de $\Lambda^i\Phi_{\bullet}$, savoir $\binom{n+1}{i}$. Donc si $L_{\bullet}$ est donné par le produit de la prop. ---.

\[
\operatorname{rg}L_n=\sum_{i\leq n+1}(\mu_i+\beta_i)\binom{n+1}{i}
+\sum_{0\leq j\leq n}\nu_j\binom{n}{i}.
\]

\uline{Donc} \unsure{1$'$)|h$'$)|2)} $L_{\bullet}$ \deleted{\unsure{en}} est \inserted{$\partial$-parfait, il est} connu à isom (non unique) près quand on connaît a) $\operatorname{rg}\pi_i(A_{\bullet})=\nu_i$ et le nombre $\beta_{i+1}$ de ses diviseurs de torsion, pour $\forall\,i\geq0$, et b) la fonction $f_{L_{\bullet}}$ \unsure{res.|res|très} pour $n$ grand (qui coïncide \deleted{alors} avec \inserted{pour $n$ grand} le polynôme $\sum_{i\geq1}(\mu_i+f_i)\binom{n+1}{i}+\sum_{j\geq0}\nu_j\binom{n}{i}$

% archive sheet number: 18
\scanpage{20}

2) Pour que $L_{\bullet}$ soit \deleted{projectif} \inserted{flasque (projectif)}, il f. et s. qu'il n'y ait que les facteurs $(\Lambda^{\unsure{j|i}}\Phi_{\ast})^{\mu_i}$ \inserted{(resp. \unsure{$=$} $(\Lambda^{0}\Psi_{\bullet})^{\nu_0}=k_{\bullet}^{\nu_0}$)}. Pour que $L_{\bullet}\otimes_k K$ (comme $K$ module $\frac12$ simplicial, où $K$ est le corps des fractions) soit flasque il f. et suffit qu'il n'y ait que des $(\Lambda^i\Phi_{\ast})^{\mu_i}$, $\Phi_{\ast}(i,n)^{f_{i,\ell}}$ \inserted{$(i\geq1)$} à l'exclusion des $(\Lambda^j\Psi_{\ast})^{\nu_j}$ $(j\geq0)$.

\deleted{Pour} 3) Pour déterminer à équiv. près la catégorie additive des k-modules $\frac12$ simpliciaux str. $\partial$ parfaits (resp str. $\partial$ parfaits) il suffit en principe de déterminer les hom mutuels entre les objets indécomposables

\[
\Lambda^i\Phi,\quad \Phi(i,n),\quad\Lambda^j\Psi_{\ast}\qquad(i\geq1,\ j\geq0)
\]

\inserted{un pour chaque $i$, et chaque diviseur $\ell$}

qu'on peut réduire à deux types si on admet aussi le cas $n=1$,

\[
\Phi(i,n)\qquad\Lambda^j\Psi_{\ast}\qquad(i\geq1,\ j\geq0)
\]

On trouve, en repassant aux complexes de chaînes associés

% archive sheet number: 19
\scanpage{21}

On \unsure{a|en|\deleted{a}} .

\[
\operatorname{Hom}_k(\Phi(i,n),L_{\bullet})\simeq d_{L_{\bullet},i}^{-1}(nL_{i-1})\subset L_i
\]

\diagram{\begin{tikzcd}[ampersand replacement=\&,column sep=small,row sep=small]
\& 1 \& n\cdot1 \&\\
0\arrow[r] \& k\arrow[u]\arrow[r,"n"]\arrow[d,"u_i"'] \& k\arrow[r,dash]\arrow[d,"u_{i-1}"] \& 0\\
0\arrow[r] \& L_i\arrow[r,"d_i"] \& L_{i-1}\arrow[r,dash] \& 0\\
\& \xi_i \& n\xi_{i-1} \&
\end{tikzcd}}

donc \deleted{$\operatorname{Hom}(\Phi(i,n),\Phi_j(m))=0$ si $j\ne\unsure{j|i},\,i+1$}

\[
\operatorname{Hom}(\Phi(i,n)_{\ast},\Phi(j,m)_{\ast})\simeq
\begin{cases}
0\ \text{si } j\ne i,\,i+1\\
\sim\{\xi\in k\mid m\deleted{\unsure{\lambda}}\xi\in nk\}=n'k\simeq k\\[-2pt]
\quad(\text{ou } n'=\frac{n}{(m,n)})\ \text{si }j=i\\
k,\ \text{si }j=i+1
\end{cases}
\]

[engendré par \unsure{la fl composée|le composé} $\Phi(i,n)\xrightarrow{\mathrm{can}}\Lambda^i\Psi_k\xrightarrow{\mathrm{can}}\Phi(i+1,m)_{\ast}$]

[p. ex. on trouve

\[
\operatorname{Hom}_k(\Phi(i,n)_{\ast},\Phi(i,\unsure{n|k})_{\ast})\simeq\deleted{\unsure{k}}k\,\mathrm{id}_{\Phi(i,n)}\simeq k
\]

(isom \deleted{de} \inserted{d'anneaux, de} $k$-algèbres si $k$ comm)

\[
\deleted{\operatorname{Hom}(\Lambda^{i}\Phi,\Lambda^{i+1}\Phi)\simeq}\qquad \forall\,i\geq0,\ n\in k-\{0\}
\]

[en particulier $\operatorname{Hom}_k(\Lambda^i\Phi_{\ast},\Lambda^i\Phi_{\ast})\simeq k\,\mathrm{id}_{\Phi(i,n)}\simeq k$ \unsure{]} 

\[
\begin{aligned}
\operatorname{Hom}_k(\Lambda^i\Phi_{\ast},\Lambda^{i+1}\Phi_{\ast})&\simeq k\cdot\lambda\simeq k,\ \text{où}\\
\lambda&=\unsure{T_{\wedge}|\top_{\wedge}}\mathpunct{:}\Lambda^i\Phi_{\ast}\longrightarrow\Lambda^{i+1}\Phi_{\ast}
\end{aligned}\quad\text{]}
\]

\[
\operatorname{Hom}_k(\Phi(i,n),\Lambda^j\Psi_{\ast})\simeq
\begin{cases}0&\text{si }j\ne i\\k&\text{si }j\simeq i\end{cases}
\]

[en particulier

\[
\operatorname{Hom}_k(\Lambda^i\Phi_{\ast},\Lambda^j\Psi_{\ast})\simeq
\begin{cases}0&\text{si }j\ne i\\k&\text{si }j=i\end{cases}
\]

et le générateur dans le cas $j=i$ est l'hom can $\Lambda^i\Phi_{\ast}\mathrel{\overset{\unsure{u^i|\alpha^i|\mathrm{can}}}{-}}\Lambda^j\Psi_{\ast}$ ]

\[
\operatorname{Hom}_k(\Lambda^i\Psi_{\ast},L_{\bullet})\simeq Z_i(\unsure{L|K|A}_{\ast})\simeq k^{\mu_{i+1}+f_{i+1}+\nu_i}
\]

en particulier

% archive sheet number: 20
\scanpage{22}

\[
\operatorname{Hom}_k(\Lambda^{\unsure{i|j}}\Psi_{\ast},\Lambda^j\Psi_{\ast})\simeq
\begin{cases}0&\text{si }j\ne i\\k&\text{si }j=i\end{cases}
\]

de \unsure{façon|\deleted{façon}} précise

\[
\operatorname{Hom}_k(\Lambda^i\Psi_{\ast},\Lambda^i\Psi_{\ast})\simeq k\,\mathrm{id}_{\Lambda^i\Psi_{\ast}}\xleftarrow{\sim} k,
\]

et

\[
\operatorname{Hom}_k(\Lambda^i\Psi_{\ast},\Phi_{\deleted{\unsure{j|\bullet}}}(j,n))\simeq
\begin{cases}0&\text{si }j\ne\deleted{\unsure{i,}}i+1\\k&\text{si }j=i+1\end{cases}
\]

en particulier

\[
\operatorname{Hom}_k(\Lambda^i\Psi_{\ast},\Lambda^j\Phi_{\ast})\simeq
\begin{cases}0&\text{si }j\ne i+1\\k&\text{si }j=i+1\end{cases}
\]

et

\[
\operatorname{Hom}_k(\Lambda^i\Psi_{\ast},\Lambda^{i+1}\Phi_{\ast})=k\cdot\unsure{i|\iota|\partial}\xleftarrow{\sim}k
\]

où $i:\Lambda^i\Psi_{\ast}\hookrightarrow\Lambda^{i+1}\Phi_{\ast}$ est l'inclusion canonique.

\uline{NB} On a un \inserted{hom de} suites exactes can

\diagram{\begin{tikzcd}[ampersand replacement=\&,column sep=small,row sep=large]
0\arrow[r,dash] \& \Lambda^{i-1}\Psi_{\ast}\arrow[r]\arrow[d,"n\,\mathrm{id}"'] \& \Lambda^{i\deleted{\unsure{+1}}}\Phi_{\ast}\arrow[r]\arrow[d,equal] \& \Lambda^i\Psi_{\ast}\arrow[r,dash]\arrow[d,Rightarrow,"\mathrm{id}"] \& 0\\
0\arrow[r,dash] \& \Lambda^{i-1}\Psi_{\ast}\arrow[r] \& \deleted{\unsure{\Lambda^i}}\Phi_{\ast}(i,n)\arrow[r,dash] \& \Lambda^i\Psi_{\ast}\arrow[r,dash] \& 0
\end{tikzcd}}

qui montre que $\Phi_{\ast}(i,n)$ est l'ext déduite de \inserted{l'ext} $\Lambda^i\Phi_{\ast}$ de $\Lambda^i\Psi_{\ast}$ par $\Lambda^{i-1}\Psi_{\ast}$ par l'hom $n\,\mathrm{id}:\Lambda^{i-1}\Psi_{\ast}\mathrel{-}\Lambda^{i-1}\Psi_{\ast}$.

% archive sheet number: 21
\scanpage{23}

\uline{En résumé} Dans la \uline{catégorie abélienne} \inserted{$k$-lin. $C$} de tous les $k$-modules $\frac12$ simpliciaux, ceux qui sont \inserted{str. $\partial$-}parfaits forment une sous-catégorie \uline{additive} $C_0$ \deleted{\unsure{engendr.}}, dans laquelle on a un complexe de cochaînes

\diagram{\begin{tikzcd}[ampersand replacement=\&,column sep=small]
0\arrow[r,dash] \& \Phi_{\ast}^{(0)}\arrow[r,dash] \& \Phi_{\ast}^{(1)}\arrow[r,dash] \& \Phi_{\ast}^{(2)}\arrow[r] \& \cdots\arrow[r] \& \Phi_{\ast}^{(i)}\arrow[r,"\lambda_i"] \& \Phi_{\ast}^{(i+1)}
\end{tikzcd}}

satisfaisant à ceci \par\textcircled{a} Complexe tot. acyclique dans $C$, \deleted{$\unsure{\lambda_i\ne0\ (i\geq0)}$}

\textcircled{b} Pour tout $i,j\geq0$, on a

\[\operatorname{Hom}(\Phi_{\ast}^{(i)},\Phi_{\ast}^{(j)})\simeq
\begin{cases}0\ \text{si }j\unsure{=|\ne}i,\,i+1\\
\simeq k\ (\text{engendré par id})\ \text{si }j=i\\
\simeq k\ (\text{----- }\lambda_i)\ \text{si }j=i+1\end{cases}\]

\deleted{\textcircled{c}} \margin{conséquences de \textcircled{a} \textcircled{b}}% slanted
Posant $\Psi_{\ast}^{(i)}=Z^{i+1}(\Phi^{\unsure{\ast|\bullet}})=\operatorname{Ker}_{C}(\Phi_{\ast}^{i+1}\xrightarrow{\lambda_{i+1}}\Phi_{\ast}^{i+2})$\inserted{$\simeq\operatorname{Coker}_{C}(\Phi^{i-1}\xrightarrow{\lambda_{i-1}}\Phi^i)$} \deleted{on a donc Hom $(\Psi^i$}

\[ \Psi^0\simeq\Phi^0 \]

\diagram{\begin{tikzcd}[ampersand replacement=\&,column sep=small,row sep=small]
(*)\quad 0\arrow[r,dash] \& \Psi^{i-1}\arrow[r,hook,"\unsure{\iota_i|v_i|\partial_i}"] \& \Phi^i\arrow[r,two heads,"u_i"]\arrow[dr,"\lambda_i"'] \& \Psi^i\arrow[r]\arrow[d,hook] \& 0 \quad\text{ext. can.}\\
\& \Phi^{i-1}\arrow[u,two heads,"u"']\arrow[ur,"\lambda_{i-1}"] \& \& \Phi^{i+1} \&
\end{tikzcd}}

\[\operatorname{Hom}(\Psi^i,\Psi^j)\ (\subset\operatorname{Hom}(\Phi^i,\Phi^{j+1}))=
\begin{cases}0&\text{si }j\ne i\deleted{\unsure{,i+1}}\\k&\text{si }j=i\ (\text{base id})\end{cases}\]

\[\operatorname{Hom}(\Phi^i,\Psi^j)\deleted{\unsure{\simeq}}[\subset\operatorname{Hom}(\Phi^i,\Phi^{j+1})]\simeq
\begin{cases}0&\text{si }j\ne i\\k&\text{si }j=\unsure{0|i|\theta}\ (\text{base }u_i)\end{cases}\]

\[\operatorname{Hom}(\Psi^{i\deleted{\unsure{+1}}},\Phi^j)\ (\subset\operatorname{Hom}(\Phi^i,\Phi^j))\simeq
\begin{cases}0&\text{si }j\ne i+1,\ (i\ne0)\\k&\text{si }j=i+1\ (\text{base }\unsure{v_i|\iota_i|\partial_i})\\ k&\text{si }j=i=0\ (\text{base id})\end{cases}\]

% slanted
\margin{\unsure{zéro car forcé|engendré} par le comp. \unsure{du hom} \unsure{canonique} $\Phi(i,n)\longrightarrow$ $\Psi^i\longrightarrow\unsure{\Phi^{i+1}|\Phi(i+1,m)}$}

d) Soit, \deleted{\unsure{$\Phi(i,n)$}} pour $n\in k-0$\inserted{et $i\geq1$}, $\Phi(i,n)$ l'ext. de $\Psi^i\deleted{\unsure{i}}$ par $\Psi^{i-1}$ déduite de l'ext. $\Phi^i$ $(*)$ par $n\,\mathrm{id}_{\Psi^{i-1}}$, de sorte qu'on a un hom de suites exactes

\diagram{\begin{tikzcd}[ampersand replacement=\&,column sep=small]
0\arrow[r] \& \Psi^{i-1}\arrow[r,"v_i"]\arrow[d,"n\,\mathrm{id}"'] \& \Phi^i\arrow[r,"u_i"]\arrow[d] \& \Psi^i\arrow[r,dash]\arrow[d,"\mathrm{id}"] \& 0\\
0\arrow[r] \& \Psi^{i-1}\arrow[r] \& \Phi(i,n)\arrow[r,dash] \& \Psi^i\arrow[r,dash] \& 0
\end{tikzcd}}

(\uline{NB}\quad $\Phi(i,1)=\Phi^i$) \inserted{On a aussi}

\diagram{\begin{tikzcd}[ampersand replacement=\&,column sep=small,row sep=small]
0\arrow[r,dash] \& \Psi^{i-1}\arrow[r,dash]\arrow[d,Rightarrow,"\mathrm{id}"'] \& \Phi(i,n)\arrow[r,dash]\arrow[d] \& \Psi^i\arrow[r]\arrow[d,"n\,\mathrm{id}"] \& 0\\
0\arrow[r,dash] \& \Psi^{i-1}\arrow[r] \& \Phi^i\arrow[r] \& \Psi^i\arrow[r] \& 0
\end{tikzcd}}

Alors on a

\[\operatorname{Hom}(\Phi(i,n),\Phi(j,m))\longrightarrow\operatorname{Hom}(\Psi^{i-1},\unsure{\Phi(j)|\Psi^j})\]

et si $j\ne i\unsure{+1|-1}$, ce dernier Hom est nul, d'où si $j\ne i\unsure{+1|-1}$, la structure d'extension est respectée par tt hom $\alpha:\Phi(i,n)\longrightarrow\Phi(j,m)$

\diagram{\begin{tikzcd}[ampersand replacement=\&,column sep=small,row sep=small]
0\arrow[r,dash] \& \Psi^{i-1}\arrow[r]\arrow[d] \& \Phi(i,n)\arrow[r]\arrow[d] \& \Psi^i\arrow[r,dash]\arrow[d] \& 0\\
0\arrow[r] \& \Psi^{j-1}\arrow[r,dash] \& \Phi(j,m)\arrow[r,dash] \& \Psi^j\arrow[r,dash] \& 0
\end{tikzcd}}

si \inserted{de plus} $j\ne i$, les flèches entre les extrêmes sont nulles, donc on trouve

\[\operatorname{Hom}(\Phi(i,n),\Phi(j,m))\xleftarrow{\sim}\operatorname{Hom}(\Psi^i,\Psi^{j-1})\quad\begin{matrix}\Vert\\0\ \text{si }j\ne i+1\end{matrix}\]

\uline{Cas $j=i+1$} \quad \textquotedblleft\quad \unsure{\textquotedblright} $\simeq k$

\uline{Cas $j=i$}. Tout hom induit une $p\,\mathrm{id}$ sur $\Psi^{i-1}$ et $q\,\mathrm{id}$ sur $\Psi^{i}$,

\diagram{\begin{tikzcd}[ampersand replacement=\&,column sep=small,row sep=small]
0\arrow[r,dash] \& \Psi^{i-1}\arrow[r,dash]\arrow[d,"p\,\mathrm{id}"'] \& \Phi(i,n)\arrow[r,dash]\arrow[d,"\alpha"] \& \Psi^i\arrow[r]\arrow[d,"q\,\mathrm{id}"] \& 0\\
0\arrow[r,dash] \& \Psi^{i-1}\arrow[r] \& \Phi(i,m)\arrow[r,dash] \& \Psi^i\arrow[r,dash] \& 0
\end{tikzcd}}

$p,q$ unique\textsuperscript{mt} déterminés par $\alpha$. On a alors (pour qu'il existe $\alpha$ définissant $p$ et $q$) que l'élément de $\operatorname{Ext}^1(\Psi^i,\Psi^{i-1})$ défini par $q^*(e_{j,m})-p_*(e_{i,n})$ (où $e_{i,n}$, $e_{j,m}$ définissent resp. les extensions $\Phi(i,n)$, $\Phi(j,m)$)

% archive sheet number: 22
\scanpage{24}

est nul ; or

\[q^*(e_{\unsure{j|i},m})=m e_i,\qquad p_*(e_{i,n})=\unsure{n e_i|mn e_i}\]

où $e_i$ définit l'extension $\Phi_i$. Donc la \unsure{condition|question} est que

\[(qm-pn)e_i=0\]

On fait l'hypothèse

\[\text{\textcircled{c}}\quad\boxed{e_i\ \text{est une base de }\operatorname{Ext}^{1}_{C_0}(\Psi^i,\Psi^{i-1})}\]

de sorte que la condition \unsure{désirée|ci-dessus|devient} est

\[ qm=pn \]

\[ \text{i.e.}\qquad \frac qp=\frac{\unsure{m|n}}{\unsure{n|m}}=\frac{m'}{n'} \]

si $\unsure{m|m'}=\frac{m}{(m,n)}$, $n'=\frac{n}{(m,n)}$, donc on aura

\[p=\deleted{\unsure{p',\,q}}\,t\unsure{n'|m'},\quad q=tm'\qquad (t\in k)\]

Or pour $p,q$ \unsure{correspondant|correspondants} ainsi à un $t$ (\unsure{évidemt|d'ailleurs} uniquement déterminé par $p,q$) l'hom $\alpha$ qui leur correspond est déterminé mod un $\deleted{\unsure{\operatorname{Ext}|\operatorname{Ex}}}\operatorname{Hom}(\Psi^i,\Psi^{i-1})=0$, donc est unique. On trouve donc

\[ \operatorname{Hom}(\Phi(i,n),\Phi(i,m))\simeq k \]

% archive sheet number: 23
\scanpage{25}

\uline{Cas} $j=i-1$

\diagram{\begin{tikzcd}[ampersand replacement=\&,column sep=small,row sep=small]
\& \Psi^{i-1}\arrow[r,dash]\arrow[d,"n\,\mathrm{id}"'] \& \Phi^{(i)}\arrow[r]\arrow[d] \& \Psi^i\arrow[d,equal] \&\\
0\arrow[r,dash] \& \Psi^{i-1}\arrow[r,dash,"D"']\arrow[d,<->,decorate,decoration={snake,amplitude=.35mm,segment length=2mm}] \& \Phi(i,n)\arrow[r,dash]\arrow[d,"\alpha=?"] \& \Psi^i\arrow[r,dash]\arrow[d,<->,decorate,decoration={snake,amplitude=.35mm,segment length=2mm}] \& 0\\
0\arrow[r] \& \Psi^{i\deleted{\unsure{-1|+2}}-2}\arrow[r] \& \Phi(i-1,m)\arrow[r,dash,"u"] \& \Psi^{i-1\deleted{\unsure{+1}}}\arrow[r,dash] \& 0
\end{tikzcd}}

le composé $\Phi^i\longrightarrow\Phi(i,n)\mathrel{-}\Phi(i-1,m)\mathrel{-}\Psi^{i-1}$ est nul car $\operatorname{Hom}(\Phi^i,\Psi^{i-1})=0$, donc l'hom \deleted{$\unsure{\Psi^i\longrightarrow\Psi|\Psi^i\longrightarrow\Psi^i}$} \inserted{$\Theta=u\alpha{D}:$} $\Psi^{i-1}\mathrel{-}\Psi^{i-1}$ composé avec $n\,\mathrm{id}$ est nul, i.e. $n\Theta=0$, or $\operatorname{Hom}(\Psi^{i-1},\Psi^{i-1})\simeq k$, donc $\Theta=0$, donc $\alpha$ respecte la structure d'extension, et on termine comme plus haut pour conclure $\alpha=0$. Donc en résumé

\[
\operatorname{Hom}(\Phi(i,n),\Phi(j,m))\simeq
\begin{cases}
0\ \text{si } j\unsure{=|\ne}i,\,i+1\\
k\ \text{si }j=i+1\\
\quad(\text{base le composé d'hom can.}\\
\quad\Phi(i,n)\mathrel{-}\Psi^i\longrightarrow\Phi(i+1,m)\unsure{)}\\
k\ \text{si }j=i\ (\text{un hom }\deleted{\text{\unsure{respecte|respectant}}}\ \text{la struct. d'ext.}
\end{cases}
\]

et étant déterminé par ce qu'il induit sur $\Psi^i$ et $\Psi^{i-1}$, qui doit être de la forme $t\,m'\cdot\mathrm{id}_{\Psi^i}$ et $t\,n'\cdot\mathrm{id}_{\Psi^{i}}$, où $n'=\frac n{(m,n)}$, $m'=\frac m{(m,n)}$ (il y a une base can. une fois qu'on a choisi le pgcd $(m,n)$ de $m$ et $n$).

% archive sheet number: 24
\scanpage{26}

\diagram{\begin{tikzcd}[ampersand replacement=\&,column sep=small,row sep=small]
\operatorname{Hom}(\Phi(i,n),\Psi^j)\arrow[ddd] \& \subset\operatorname{Hom}(\Phi(i,n),\Phi^{j+1})\\
\& \begin{matrix}\Vert\\\Phi(j+1,\unsure{n|1|-})\end{matrix}\\
=0\ \text{si} \& j+1\ne i,i+1\quad\text{i.e. si }j\ne i-1,i\\
\operatorname{Hom}(\Psi^{i-1},\Psi^j)=0\ \text{si }j\ne i-1 \&
\end{tikzcd}}

Si donc $j=i$, $\operatorname{Hom}(\Phi(i,n),\Psi^{j=i})\xleftarrow{\sim}\operatorname{Hom}(\unsure{\Psi^i|\Phi^i},\Psi^i)$ $\underset{k}{\unsure{\simeq|=}}$ (base l'hom can)

\deleted{Si} Si $j=i-1$,

\diagram{\begin{tikzcd}[ampersand replacement=\&,column sep=small,row sep=small]
{}-{}\& \Psi^{i-1}\arrow[r]\arrow[d,"n\,\mathrm{id}"'] \& \Phi^i\deleted{\unsure{(1)|(k)}}\arrow[r]\arrow[d] \& \Psi^i\arrow[r,dash] \& 0\\
0\arrow[r,dash] \& \Psi^{i-1}\arrow[r,"\unsure{\iota|i|v}"] \& \Phi(i,n)\arrow[r,dash]\arrow[d,"\alpha=?"] \& \Psi^i\arrow[r,dash]\arrow[u] \& 0\\
\& \& \Psi^{i-1} \& \&
\end{tikzcd}}

le composé de $\alpha\deleted{\unsure{\iota|i}}$ avec \deleted{$\unsure{n\,\mathrm{id}_{\Psi^i}|u\circ d_{\Psi^i}}$} \inserted{$\Phi^i\mathrel{-}\Phi(i,n)$} est $\in\operatorname{Hom}(\Phi^i,\Psi^{i-1})=0$, donc \deleted{$n\alpha\unsure{\iota|i}\in\operatorname{Hom}$} \inserted{son composé} $\alpha\iota\circ n\,\mathrm{id}_{\Psi^{i-1}}=n(\alpha\iota)\in\operatorname{Hom}(\Psi^{i-1},\Psi^{i-1})$ est nul, et comme $\operatorname{Hom}(\Psi^{i-1},\Psi^{i-1})\simeq k$, on voit que $\alpha\iota=0$, donc $\alpha$ \deleted{\unsure{passe par|provient}} se factorise par $\Psi^i\mathrel{-}\Psi^{i-1}$, donc est nul. Donc

\[\operatorname{Hom}(\Phi(i,n),\Psi^j)=
\begin{cases}0&\text{si }j\ne i\\
\simeq k\ (\text{base l'hom can }\Phi(i,n)\mathrel{-}\Psi^i)&\text{si }j=i\end{cases}\]

% archive sheet number: 25
\scanpage{27}

\[\operatorname{Hom}(\Psi^i_{\unsure{\ast|\bullet}},\Phi(j,n))\longrightarrow
\underset{\substack{\Vert\\0\ \text{si }j\ne\unsure{j|i}}}{\operatorname{Hom}(\Psi^i,\Psi^j)}\]

donc si $j\ne i$, \inserted{$\forall\,\alpha\in$} $\operatorname{Hom}(\Psi^i_{\deleted{\unsure{\ast}}},\Phi(j,n))\xleftarrow{\sim}\operatorname{Hom}(\Psi^i,\Psi^{j-1}\deleted{\unsure{\ne i}})$, \deleted{\unsure{d'où}} qui est nul si $i\ne j-1$ i.e. $j\ne i+1$ ; si $j=i+1$, on trouve $\simeq k$ \unsure{---}.

Pour $j=i$, le composé $\Psi^i\xrightarrow{\alpha}\Phi(i,n)\xrightarrow{\mathrm{can}}\Psi^i$

\diagram{\begin{tikzcd}[ampersand replacement=\&,column sep=small,row sep=small]
\& \& \Psi^i\arrow[d,"\alpha"']\arrow[dr] \& \&\\
0\arrow[r] \& \Psi^{j-1}\arrow[r,dash] \& \Phi(j,n)\arrow[r,dash] \& \Psi^j\arrow[r,dash] \& 0
\end{tikzcd}}

est de la forme $d\cdot\mathrm{id}_{\Psi^i}$, donc l'image inverse \deleted{\unsure{par|p}}de l'extension $\Phi(i,n)$ est de classe $d e_{i,n}=dn e_i$, qui est nulle, donc $dn e_i=0$, donc $dn=0$, donc $d=0$, donc \deleted{\unsure{$\Phi$}} $\alpha$ se factorise par $\Psi^i\mathrel{-}\Psi^{i-1}$, donc est nul.

\deleted{Pour $j=i+1$, on a $\operatorname{Hom}($} Donc \deleted{$\unsure{\operatorname{Hom}(\Phi,\Psi}$}

\[\operatorname{Hom}(\Psi^i,\Phi(j,n))\simeq
\begin{cases}0\ \text{si }j\ne i+1\\ k\ \text{si }j=i+1\ (\text{base \unsure{formé|formée} de}\\
\quad\text{l'hom can }\Psi^i\mathrel{-}\Phi(i+1,n))\end{cases}\]

\noindent\rule{.64\linewidth}{.3pt}

Enfin on suppose

\begin{tabular}{|p{.85\linewidth}}\textcircled{d} Tout objet de $C$ est somme finie d'objets de la forme $\Phi(i,n)$ \inserted{$(\unsure{i|n}\geq1)$} ou $\Psi^j$ $(j\geq0)$\end{tabular}

% archive sheet number: 26
\scanpage{28}

En vertu \unsure{du critère|des axiomes|on sait} qui \unsure{suit}, si $A$ est la sous-catégorie pleine de $C$ engendrée par les $\Phi(i,n)$ et les $\Psi^i$, alors $C$ est connue à équiv près (\unsure{déf à isom.|déf. isom.} unique près) si on connaît $A$ avec sa structure $\frac12$-additive. Or c'est \inserted{elle} qui vient d'être déterminée (on a omis de décrire la composition des morphismes, qui est évidente \ldots).

On trouve donc que $C\simeq$ groupes $\frac12$ simpliciaux abéliens str. $\partial$-parfaits \ldots.

La chose restant à faire est d'en déterminer la $\otimes$-structure et les opérations $\Gamma^{\unsure{i|\prime}}$.

% archive sheet number: 27
\scanpage{29}

Soit $C$ une catégorie, $A$ une sous-catégorie \deleted{\unsure{Con}} \inserted{pleine}. Considérons \unsure{le foncteur composé|les foncteurs composés}

\diagram{\begin{tikzcd}[ampersand replacement=\&,column sep=large]
C\arrow[r,"\mathrm{can}"]\arrow[rr,bend right=22,"\unsure{\rho|r}"'] \& \widehat C\arrow[r,"\mathrm{rest.}"] \& \widehat A=\underline{\operatorname{Hom}}(A^\circ,\mathrm{Ens})
\end{tikzcd}}

si \deleted{tout objet $X\in C$} tout objet \inserted{$X$} de $C$ est $\varinjlim$ \inserted{dans $C$} d'un diagramme de $A$ (ou ce qui revient au $=$, $X\simeq\varinjlim_{Y\in\mathrm{Ob}\,A/X}Y$), alors le foncteur $\rho$ est pl. fidèle.

Si $C$ est additive \inserted{(mais pas néc $A$)}, le diagramme ci-dessus se complète en

\diagram{\begin{tikzcd}[ampersand replacement=\&,column sep=large,row sep=large]
C\arrow[r,"c_{\mathrm{ab}}"]\arrow[rr,bend left=27,"r_{\mathrm{ab}}"]\arrow[dr,"c"']\arrow[drr,bend right=32,"r"'] \& \widehat C_{\mathrm{ab}}\arrow[r,"\mathrm{rest.}\,\rho_{\mathrm{ab}}"]\arrow[d,"\mathrm{oubli}\ \omega_C"] \& \widehat A_{\mathrm{ab}}\arrow[d,"\mathrm{oubli}\ \omega_A"]\\
\& \widehat C\arrow[r,"\mathrm{rest}\,\rho"] \& \widehat A
\end{tikzcd}}

% slanted
\margin{\unsure{itou|idem|i.e.} \unsure{p\textsuperscript{r}|pour} structure $k$-lin. sur $C$}

avec $\omega_A$ fidèle, d'où il résulte que $r_{\mathrm{ab}}$ est aussi pl. fidèle.

[C'est le cas en particulier si $\forall\,X\in C$ est $\simeq$ somme d'objets de $C$]

Supposons de plus que tout $X\in C$ soit aussi \unsure{isom.|somme} à une $\varprojlim$ \inserted{dans $C$} d'objets de $C$. Alors l'image essentielle de $C$ dans $\widehat A$ (resp. $\widehat A_{\mathrm{ab}}$) est contenue dans la sous cat. pleine de $\widehat A$ resp $\widehat A_{\mathrm{ab}}$ formée des foncteurs

% archive sheet number: 28
\scanpage{30}

qui sont $\varprojlim$ du type envisagé de foncteurs représentables (dans le cas resp., on utilise, pour \deleted{\unsure{y}} donner un sens à cet énoncé -- dans $\widehat A_{\mathrm{ab}}$ -- la structure \textquotedblleft$\frac12$ additive\textquotedblright\ sur $A$ héritée par son plongement dans $C$ additive). Lorsque de plus $C$ est stable pour ce type de $\varprojlim$ (p. ex. des produits finis), alors on a \unsure{$=$|même} une \uline{équivalence} entre $C$ et la cat. des foncteurs $A^\circ\mathrel{-}\mathrm{Ens}$ (resp $A^\circ\mathrel{-}\mathrm{Ab}$) qui sont de telles $\varprojlim$ de foncteurs représentables.

\deleted{[Si par ex. $A$ est réduite à un objet, on trouve dans le cas $C$ additive] \unsure{cas resp., si tt objet|resp., si $A$ objet}}

Si p. ex. $C$ est additive, donc \deleted{$A$} induit sur $A$ une structure $\frac12$ additive, et si tt objet de $C$ est isomorphe à une somme directe \inserted{finie} d'objets de $A$, alors $C\simeq$ sous-cat. pleine de $\widehat A_{\mathrm{ab}}$ (foncteurs additifs contravariants de $A$ dans $\mathrm{Ab}$) formée \deleted{sont} des foncteurs isom. à une somme directe finie de foncteurs repr.

% archive sheet number: 29
